\documentclass[12pt,a4paper]{article}

\usepackage[utf8]{inputenc}
\usepackage[english,russian]{babel}
\usepackage{amsmath,amsfonts,amssymb,amsthm}
\usepackage{graphicx}
\usepackage{hyperref}
\usepackage{geometry}
\usepackage{algorithm}
\usepackage{algorithmic}

\geometry{margin=2.5cm}

% --- Theorem environments ---
\newtheorem{theorem}{Теорема}
\newtheorem{lemma}{Лемма}
\newtheorem{proposition}{Утверждение}
\newtheorem{corollary}{Следствие}

\theoremstyle{definition}
\newtheorem{definition}{Определение}
\newtheorem{remark}{Замечание}

% -------------------------------------------------

\title{Ловушки дефекции в многоагентном Q-обучении:\\
Стохастическая теория и численная верификация}
\author{Проект RL\_GameTheory}
\date{\today}

\begin{document}

\maketitle

% =================================================
\begin{abstract}
Рассматривается феномен \emph{ловушек дефекции} в многоагентном Q-обучении
с Boltzmann-политикой в обобщённой дилемме заключённого.
Показано, что при малом шаге обучения $\alpha$ система демонстрирует
метастабильное дефектное равновесие, выход из которого возможен
только за счёт стохастических флуктуаций.

Получена асимптотическая формула для среднего времени выхода из ловушки
в пределе $\alpha \to 0$:
\[
\mathbb{E}[\tau_{\mathrm{exit}}]
\sim
\exp\!\left(
\frac{\kappa \Delta^2}{\alpha \Sigma^2(x^*)}
\right),
\]
где $\kappa$ — локальная сила возврата к равновесию,
$\Delta$ — расстояние до границы ловушки,
$\Sigma^2(x^*)$ — дисперсия стохастического шума в равновесии.
\end{abstract}

% =================================================
\section{Введение}

Многоагентное обучение с подкреплением в играх приводит
к сложной стохастической динамике, включающей
метастабильные состояния и чрезвычайно медленную сходимость.
Одним из наиболее устойчивых феноменов являются
\emph{ловушки дефекции} — длительное пребывание системы
в состоянии почти полной дефекции,
несмотря на наличие более выгодных кооперативных исходов.

В данной работе мы даём строгую стохастическую теорию
возникновения и разрушения таких ловушек
в многоагентном Q-обучении с Boltzmann-политикой
и выводим экспоненциальную асимптотику
времени выхода из ловушки.

% =================================================
\section{Модель}

\subsection{Обобщённая дилемма заключённого}

Рассматривается $n$-игроковая линейная игра общественных благ.

\begin{definition}
Каждый игрок выбирает действие $a_i \in \{C,D\}$.
Пусть
\[
k = \sum_{i=1}^n \mathbf{1}_{\{a_i=C\}}
\]
— число кооператоров.
Выплата игрока $i$ равна
\[
u_i =
\begin{cases}
\dfrac{b k}{n} - c, & a_i = C,\\[0.4em]
\dfrac{b k}{n}, & a_i = D,
\end{cases}
\]
где $b>0$ — выгода от кооперации, $c>0$ — её стоимость.
\end{definition}

\begin{proposition}
Игра является дилеммой заключённого при условии
\[
\frac{b}{n} < c < b.
\]
\end{proposition}

% -------------------------------------------------
\subsection{Q-обучение с Boltzmann-политикой}

Каждый агент поддерживает значения
$Q_t(C), Q_t(D)$ и выбирает действия согласно
Boltzmann-политике
\[
\pi_t(a)
=
\frac{e^{\beta Q_t(a)}}
{e^{\beta Q_t(C)} + e^{\beta Q_t(D)}},
\qquad \beta>0.
\]

Обновление Q-значений имеет вид
\[
Q_{t+1}(a) =
\begin{cases}
Q_t(a) + \alpha
\bigl(
r_t + \gamma \max_{a'} Q_t(a') - Q_t(a)
\bigr),
& a=a_t,\\
Q_t(a), & \text{иначе}.
\end{cases}
\]

% =================================================
% =================================================
\section{Динамика}

В данном разделе мы строго выводим стохастическую динамику,
описывающую эволюцию многоагентного Q-обучения.

% -------------------------------------------------
\subsection{Редукция к разности Q-значений}

\begin{lemma}[Редукция к разности Q-значений]
\label{lem:reduction}
В симметричном режиме многоагентного Q-обучения
динамика системы полностью описывается
одномерным процессом
\[
\Delta Q_t := Q_t(C) - Q_t(D).
\]
Вероятность выбора действия $C$ в момент времени $t$ равна
\[
p_t = \sigma(\beta \Delta Q_t),
\qquad
\sigma(x) = \frac{1}{1+e^{-x}}.
\]
\end{lemma}

\begin{proof}
По определению Boltzmann-политики,
вероятность выбора действия $C$ равна
\[
\pi_t(C)
=
\frac{e^{\beta Q_t(C)}}
{e^{\beta Q_t(C)} + e^{\beta Q_t(D)}}.
\]
Вынесем множитель $e^{\beta Q_t(D)}$ из знаменателя:
\[
\pi_t(C)
=
\frac{e^{\beta (Q_t(C)-Q_t(D))}}
{1 + e^{\beta (Q_t(C)-Q_t(D))}}
=
\sigma(\beta \Delta Q_t).
\]

Таким образом, стратегия агента зависит от Q-значений
только через их разность $\Delta Q_t$.
Абсолютный сдвиг
$Q_t(C)\mapsto Q_t(C)+K$, $Q_t(D)\mapsto Q_t(D)+K$
не влияет ни на стратегию, ни на будущие обновления.

Так как все агенты симметричны и рассматривается
симметричное подпространство состояний,
эволюция системы редуцируется
к динамике одной скалярной величины $\Delta Q_t$.
\end{proof}

% -------------------------------------------------
\subsection{Стохастическая аппроксимация}

Обозначим через $\mathcal{F}_t$
состояние системы, порождённое процессом Q-обучения,
то есть $\mathcal{F}_t$ содержит всю информацию
о действиях, наградах и Q-значениях до момента $t$ включительно.

\begin{lemma}[Стохастическая аппроксимация]
\label{lem:sa}
Существует измеримая функция $F:\mathbb{R}\to\mathbb{R}$
и последовательность случайных величин $\{\eta_t\}_{t\ge 0}$,
такие что
\[
\Delta Q_{t+1}
=
\Delta Q_t + \alpha F(\Delta Q_t) + \alpha \eta_t,
\]
где выполнены условия:
\begin{align*}
\mathbb{E}[\eta_t \mid \mathcal{F}_t] &= 0,\\
\sup_{t} \mathbb{E}[\eta_t^2 \mid \mathcal{F}_t] &< \infty.
\end{align*}
\end{lemma}

\begin{proof}
По правилу Q-обновления имеем:
\[
Q_{t+1}(a)
=
Q_t(a)
+
\alpha \mathbf{1}_{\{a_t = a\}}
\bigl(
r_t + \gamma \max_{a'} Q_t(a') - Q_t(a)
\bigr),
\quad a\in\{C,D\}.
\]

Вычитая обновление для $D$ из обновления для $C$, получаем:
\[
\Delta Q_{t+1} - \Delta Q_t
=
\alpha \xi_t,
\]
где случайная величина $\xi_t$ равна
\[
\xi_t :=
\mathbf{1}_{\{a_t = C\}}
\bigl(
r_t + \gamma \max_{a'} Q_t(a') - Q_t(C)
\bigr)
-
\mathbf{1}_{\{a_t = D\}}
\bigl(
r_t + \gamma \max_{a'} Q_t(a') - Q_t(D)
\bigr).
\]

Определим детерминированный дрейф как условное математическое ожидание:
\[
F(x)
:=
\mathbb{E}[\xi_t \mid \Delta Q_t = x].
\]
Далее введём шумовую составляющую:
\[
\eta_t := \xi_t - F(\Delta Q_t).
\]

По построению имеем:
\[
\mathbb{E}[\eta_t \mid \mathcal{F}_t]
=
\mathbb{E}[\xi_t \mid \mathcal{F}_t]
-
F(\Delta Q_t)
=
0,
\]
то есть $\eta_t$ является мартингальной разностью
относительно фильтрации $\{\mathcal{F}_t\}$.

Так как награды $r_t$ по предположению имеют конечную дисперсию,
а Q-значения остаются конечными при фиксированном шаге обучения,
существует константа $C<\infty$ такая, что
\[
\mathbb{E}[\eta_t^2 \mid \mathcal{F}_t] \le C
\quad \text{для всех } t.
\]

Следовательно, процесс $\Delta Q_t$
является стохастической аппроксимацией
с шагом $\alpha$, детерминированным дрейфом $F$
и мартингальным шумом $\eta_t$.
\end{proof}

\begin{remark}
При формальном пределе $\eta_t \equiv 0$
получается детерминированное рекуррентное соотношение,
которое сходится к устойчивому дефектному равновесию.
Таким образом, существование ловушки дефекции
является исключительно стохастическим эффектом.
\end{remark}


% ====================================================



% =================================================
% =================================================
% =================================================
% =================================================
% =================================================
% =================================================
\section{Детерминированный дрейф в дефектной области}

В этом разделе мы явно вычисляем функцию дрейфа $F(x)$,
определённую как условное математическое ожидание
приращения разности Q-значений,
в области $\Delta Q < 0$.
Все вычисления являются точными и не используют асимптотических
приближений.

% -------------------------------------------------
\subsection{Выигрыши и TD-ошибки}

Рассмотрим фиксированного агента $i$.
Пусть $K_{-i}$ — случайная величина, равная числу кооператоров
среди остальных $n-1$ агентов.

При фиксированном значении $k$ выигрыши агента $i$ равны:
\[
r_C(k) = \frac{b(k+1)}{n} - c + \mathrm{offset},
\qquad
r_D(k) = \frac{bk}{n} + \mathrm{offset}.
\]

Для действий $a\in\{C,D\}$ определим TD-ошибки:
\[
TD_a
=
r_a(K_{-i})
+
\gamma \max(Q^C,Q^D)
-
Q^a.
\]

% -------------------------------------------------
\subsection{Инкремент разности Q-значений}

Определим разность Q-значений:
\[
\Delta Q(m) := Q^C(m)-Q^D(m).
\]

Определим нормированный инкремент:
\[
\xi(m)
:=
\frac{\Delta Q(m+1)-\Delta Q(m)}{\alpha}.
\]

\begin{lemma}[Точное выражение для инкремента]
Выполняется:
\[
\xi(m)
=
\begin{cases}
TD_C, & \text{если } a_i(m)=C,\\[0.3em]
-\,TD_D, & \text{если } a_i(m)=D.
\end{cases}
\]
\end{lemma}

\begin{proof}
Если $a_i(m)=C$, то обновляется только $Q^C$:
\[
\Delta Q(m+1)-\Delta Q(m)
=
\alpha TD_C.
\]
Если $a_i(m)=D$, то обновляется только $Q^D$:
\[
\Delta Q(m+1)-\Delta Q(m)
=
-\alpha TD_D.
\]
Деление на $\alpha$ даёт требуемый результат.
\end{proof}

% -------------------------------------------------
\subsection{Определение дрейфа}

Определим функцию дрейфа как
\[
F(x)
:=
\mathbb E\!\left[
\xi(m)\mid \Delta Q(m)=x
\right].
\]

По закону полного ожидания:
\[
F(x)
=
\sigma(\beta x)\,\mathbb E[TD_C\mid \Delta Q=x]
-
(1-\sigma(\beta x))\,\mathbb E[TD_D\mid \Delta Q=x].
\]

% -------------------------------------------------
\subsection{Средняя разность выигрышей}

\begin{lemma}[Средняя разность выигрышей]
Выполняется:
\[
\mathbb E\!\left[r_C(K_{-i})-r_D(K_{-i})\right]
=
\frac{b}{n}-c.
\]
\end{lemma}

\begin{proof}
Для любого $k$:
\[
r_C(k)-r_D(k)
=
\left(\frac{b(k+1)}{n}-c\right)-\frac{bk}{n}
=
\frac{b}{n}-c.
\]
Следовательно, это равенство сохраняется и после взятия
математического ожидания по распределению $K_{-i}$.
\end{proof}

% -------------------------------------------------
\subsection{Дрейф в дефектной области и разложение на главный вклад и остаток}

Рассмотрим шаг алгоритма $m$ и зафиксируем значение
\[
x := \Delta Q(m) = Q^C(m) - Q^D(m).
\]

Мы будем работать в дефектной области
\[
x<0,
\qquad
Q^D(m) > Q^C(m),
\]
в которой выполняется
\[
\max\{Q^C(m),Q^D(m)\} = Q^D(m).
\]

% -------------------------------------------------
\subsubsection{TD-ошибки}

Для действий $a\in\{C,D\}$ TD-ошибки имеют вид
\[
TD_a
=
r_a(K_{-i})
+
\gamma Q^D(m)
-
Q^a(m),
\]
где $K_{-i}$ — число кооператоров среди остальных агентов.

Следовательно,
\begin{align}
TD_C
&=
r_C
+
\gamma Q^D
-
Q^C,
\label{eq:TD_C}\\
TD_D
&=
r_D
+
(\gamma-1)Q^D.
\label{eq:TD_D}
\end{align}

% -------------------------------------------------
\subsubsection{Определение дрейфа}

Определим шумовую переменную
\[
\xi(m)
=
\begin{cases}
TD_C, & a_i(m)=C,\\
-\,TD_D, & a_i(m)=D.
\end{cases}
\]

Детерминированный дрейф задаётся как
\[
F(x)
:=
\mathbb E\!\left[\xi(m)\mid \Delta Q(m)=x\right].
\]

\begin{lemma}[Закон полного ожидания по действию]
\label{lem:LOTE}
При фиксированном значении $\Delta Q(m)=x$ выполняется
\begin{align}
F(x)
&=
\sigma(\beta x)\,
\mathbb E\!\left[TD_C\mid x\right]
-
\bigl(1-\sigma(\beta x)\bigr)\,
\mathbb E\!\left[TD_D\mid x\right].
\label{eq:F_LOTE}
\end{align}
\end{lemma}

\begin{proof}
Условно на $\Delta Q(m)=x$ действие агента $a_i(m)$ принимает значения
$C$ и $D$ с вероятностями $\sigma(\beta x)$ и $1-\sigma(\beta x)$
соответственно.
По закону полного математического ожидания по разбиению
$\{a_i(m)=C\}\cup\{a_i(m)=D\}$ получаем \eqref{eq:F_LOTE}.
\end{proof}

% -------------------------------------------------
\subsubsection{Точное выражение для дрейфа}

Подставляя \eqref{eq:TD_C}--\eqref{eq:TD_D} в \eqref{eq:F_LOTE}, получаем
\begin{align}
F(x)
&=
\sigma(\beta x)
\bigl(
\mathbb E[r_C]
+
\gamma Q^D
-
Q^C
\bigr)
-
\bigl(1-\sigma(\beta x)\bigr)
\bigl(
\mathbb E[r_D]
+
(\gamma-1)Q^D
\bigr).
\label{eq:F_step1}
\end{align}

Используя тождество $Q^C = Q^D + x$, справедливое при условии
$\Delta Q(m)=x$, преобразуем \eqref{eq:F_step1} к виду
\begin{align}
F(x)
&=
\sigma(\beta x)\bigl(\mathbb E[r_C]+\mathbb E[r_D]\bigr)
-
\mathbb E[r_D]
-
\sigma(\beta x)\,x
+
(\gamma-1)Q^D\bigl(2\sigma(\beta x)-1\bigr).
\label{eq:F_exact}
\end{align}

Формула \eqref{eq:F_exact} является \emph{точным} выражением для дрейфа
в дефектной области.

% -------------------------------------------------
\subsubsection{Выделение главного вклада}

В игре общественных благ выполняется
\[
\mathbb E[r_C-r_D]=\frac{b}{n}-c,
\]
и потому
\[
\mathbb E[r_C]+\mathbb E[r_D]
=
\left(\frac{b}{n}-c\right)
+
2\mathbb E[r_D].
\]

Определим главный вклад дрейфа как
\begin{equation}
F_{\mathrm{main}}(x)
:=
\sigma(\beta x)
\left(
\frac{b}{n}-c-x
\right).
\label{eq:F_main}
\end{equation}

Оставшуюся часть определим как остаток:
\[
R(x)
:=
F(x)-F_{\mathrm{main}}(x).
\]

Явно из \eqref{eq:F_exact} и \eqref{eq:F_main} получаем
\begin{align}
R(x)
&=
(\gamma-1)Q^D\bigl(2\sigma(\beta x)-1\bigr)
+
2\sigma(\beta x)\mathbb E[r_D]
-
\mathbb E[r_D].
\label{eq:R_exact}
\end{align}

% -------------------------------------------------
\subsubsection{Оценка остатка}

\begin{lemma}[Ограниченность остатка]
\label{lem:R_bound}
Пусть параметры $\beta,\gamma$ фиксированы, а выигрыши имеют конечное
математическое ожидание.
Тогда существует константа $C<\infty$, такая что
\[
|R(x)| \le C
\qquad
\text{для всех } x<0.
\]
\end{lemma}

\begin{proof}
Функция $\sigma(\beta x)$ принимает значения в интервале $(0,1)$.
Следовательно,
\[
|2\sigma(\beta x)-1|\le 1.
\]
Так как $Q^D(m)$ фиксировано на шаге $m$, а $\mathbb E[r_D]$ конечно,
каждое слагаемое в \eqref{eq:R_exact} ограничено по модулю.
Сумма ограниченных слагаемых также ограничена.
\end{proof}

\begin{remark}
Остаток $R(x)$ не содержит множителей порядка $1/\alpha$ и не масштабируется
при $\alpha\to0$.
Поэтому при анализе асимптотики времени выхода из ловушки
он не влияет на ведущий экспоненциальный порядок.
\end{remark}


% -------------------------------------------------
\subsection{Равновесие}

\begin{theorem}[Дефектное равновесие]
\label{thm:defection_equilibrium}
Если $\frac{b}{n}<c$, то уравнение
\[
F_{\mathrm{main}}(x)=0
\]
имеет единственное решение
\[
x^\star=\frac{b}{n}-c<0.
\]
\end{theorem}

\begin{proof}
Так как $\sigma(\beta x)>0$ для всех $x$, равенство
$F_{\mathrm{main}}(x)=0$ эквивалентно
\[
\frac{b}{n}-c-x=0,
\]
откуда следует $x^\star=\frac{b}{n}-c$.
\end{proof}

\begin{definition}[Сила возврата]
Определим
\[
\kappa := -F_{\mathrm{main}}'(x^\star)
=
\sigma(\beta x^\star).
\]
\end{definition}



% =================================================
% =================================================
\section{Непрерывное приближение динамики}

В данном разделе мы показываем, что при малом шаге обучения $\alpha$
дискретная стохастическая динамика разности Q-значений
\[
\Delta Q(m)
\]
может быть аппроксимирована детерминированной непрерывной динамикой
на временных масштабах порядка $1/\alpha$.
Полученное приближение описывает усреднённое поведение системы
на конечных временных интервалах.

% -------------------------------------------------
\subsection{Общий вид дискретной динамики}

Рассмотрим одномерный дискретный процесс вида
\begin{equation}
\label{eq:sa_general_simple}
X_{m+1}
=
X_m
+
\alpha F(X_m)
+
\alpha \eta(m),
\end{equation}
где:
\begin{itemize}
\item $F:\mathbb{R}\to\mathbb{R}$ — детерминированная функция;
\item $\{\eta(m)\}$ — последовательность случайных величин,
удовлетворяющая условию нулевого условного среднего:
\[
\mathbb E[\eta(m)\mid \mathcal F_m]=0;
\]
\item условные вторые моменты равномерно ограничены:
\[
\sup_m \mathbb E[\eta(m)^2\mid \mathcal F_m] < \infty.
\]
\end{itemize}

В предыдущем разделе было показано, что процесс
\[
X_m=\Delta Q(m)
\]
имеет именно такой вид,
где функция $F$ задаётся формулой \eqref{eq:F_exact},
а $\eta(m)$ представляет собой центрированный стохастический шум.

\medskip

\textbf{Замечание.}
В силу структуры модели функция $F$ является непрерывной и локально
липшицевой на любом компакте, содержащем значения $\Delta Q(m)$
на конечном временном интервале.
Это следует из гладкости сигмоидальной функции $\sigma(\beta x)$
и из того, что условные математические ожидания наград
$\mathbb E[r_C]$, $\mathbb E[r_D]$
являются гладкими функциями вероятности кооперации,
а значит, гладкими функциями $x=\Delta Q$.

% -------------------------------------------------
\subsection{Масштабирование времени}

Для анализа поведения процесса на больших временах
введём новое непрерывное время
\[
s:=\alpha m
\]
и определим интерполированный процесс
\[
X^{(\alpha)}(s)
:=
\Delta Q\bigl(\lfloor s/\alpha \rfloor\bigr).
\]

Таким образом, $X^{(\alpha)}(s)$ является кусочно-постоянной функцией,
которая описывает дискретную динамику $\Delta Q(m)$
на временной шкале порядка $1/\alpha$.

% -------------------------------------------------
\subsection{Предел при исчезающем шаге обучения}

\begin{theorem}[Детерминированный предельный процесс]
\label{thm:ode_limit_simple}
Пусть $\alpha\to0$.
Тогда для любого $T>0$ существует функция $x:[0,T]\to\mathbb R$
такая, что процесс $X^{(\alpha)}(s)$
сходится к $x(s)$ по вероятности,
равномерно по $s\in[0,T]$,
и функция $x(s)$ удовлетворяет интегральному уравнению
\[
x(t)-x(s)
=
\int_s^t F\bigl(x(u)\bigr)\,du,
\qquad
0\le s<t\le T.
\]
В частности, функция $x$ дифференцируема и является решением
обыкновенного дифференциального уравнения
\[
\frac{dx}{ds}=F(x),
\qquad
x(0)=\lim_{\alpha\to0}X^{(\alpha)}(0).
\]
\end{theorem}

\begin{proof}
Зафиксируем $s\ge0$ и $\Delta s>0$.
Рассмотрим приращение интерполированного процесса:
\begin{align*}
X^{(\alpha)}(s+\Delta s)-X^{(\alpha)}(s)
&=
\sum_{k=\lfloor s/\alpha \rfloor}^{\lfloor (s+\Delta s)/\alpha \rfloor-1}
\alpha F(X_k)
+
\sum_{k=\lfloor s/\alpha \rfloor}^{\lfloor (s+\Delta s)/\alpha \rfloor-1}
\alpha \eta(k) \\
&=: S_1^{(\alpha)}(s,\Delta s) + S_2^{(\alpha)}(s,\Delta s).
\end{align*}

\medskip
\noindent
\textbf{Детерминированный член.}
Заметим, что число слагаемых в сумме $S_1^{(\alpha)}$
равно
\[
N_\alpha
=
\lfloor (s+\Delta s)/\alpha \rfloor-\lfloor s/\alpha \rfloor
=
\frac{\Delta s}{\alpha}+O(1).
\]
Для каждого $k$ положим $u_k:=\alpha k$.
Тогда $u_k\in[s-\alpha,s+\Delta s]$,
и соответствующий вклад равен
\[
\alpha F(X_k)=\alpha F\bigl(X^{(\alpha)}(u_k)\bigr).
\]

Следовательно, сумма $S_1^{(\alpha)}$ является римановой суммой
для интеграла
\[
\int_s^{s+\Delta s} F\bigl(X^{(\alpha)}(u)\bigr)\,du,
\]
с погрешностью порядка $O(\alpha)$,
связанной с округлением границ интервала
и кусочно-постоянной интерполяцией.
Эта погрешность стремится к нулю при $\alpha\to0$,
равномерно по $s$ на конечных интервалах.

\medskip
\noindent
\textbf{Случайный член.}
Рассмотрим сумму
\[
S_2^{(\alpha)}(s,\Delta s)
=
\sum_{k=\lfloor s/\alpha \rfloor}^{\lfloor (s+\Delta s)/\alpha \rfloor-1}
\alpha \eta(k).
\]
Используя свойство нулевого условного среднего
и итеративное применение свойства условного ожидания, получаем
\[
\mathbb E\!\left[S_2^{(\alpha)}(s,\Delta s)\mid \mathcal F_{\lfloor s/\alpha \rfloor}\right]=0.
\]

Для условной дисперсии имеем:
\[
\mathbb E\!\left[\bigl(S_2^{(\alpha)}(s,\Delta s)\bigr)^2
\mid \mathcal F_{\lfloor s/\alpha \rfloor}\right]
=
\sum_k \alpha^2
\mathbb E[\eta(k)^2\mid \mathcal F_k],
\]
где перекрёстные члены исчезают,
поскольку $\eta(k)$ имеет нулевое условное среднее
относительно $\mathcal F_k$.
Так как число слагаемых порядка $\Delta s/\alpha$,
а условные дисперсии равномерно ограничены,
получаем оценку
\[
\mathbb E\!\left[\bigl(S_2^{(\alpha)}(s,\Delta s)\bigr)^2
\mid \mathcal F_{\lfloor s/\alpha \rfloor}\right]
\le
C\,\alpha\,\Delta s
\;\xrightarrow[\alpha\to0]{}\;0.
\]
Следовательно, $S_2^{(\alpha)}(s,\Delta s)\to0$
по вероятности, равномерно на компактных интервалах.

\medskip
\noindent
\textbf{Предел.}
Из предыдущих шагов следует, что для любого $0\le s<t\le T$
приращение интерполированного процесса можно представить в виде
\[
X^{(\alpha)}(t)-X^{(\alpha)}(s)
=
\int_s^t F\bigl(X^{(\alpha)}(u)\bigr)\,du
+
R^{(\alpha)}(s,t),
\]
где остаточный член $R^{(\alpha)}(s,t)$ сходится к нулю по вероятности
при $\alpha\to0$, равномерно по всем $s,t\in[0,T]$.

Кроме того, из построения следует, что семейство функций
$\{X^{(\alpha)}\}_{\alpha>0}$
является равномерно ограниченным на $[0,T]$.
В самом деле, при конечном времени $T$
число шагов дискретной динамики пропорционально $T/\alpha$,
а каждое приращение имеет порядок $\alpha$,
так что значения $X^{(\alpha)}(s)$ остаются в компактном множестве
с вероятностью, стремящейся к единице при $\alpha\to0$.

Далее, из представления приращений следует,
что семейство $\{X^{(\alpha)}\}$
является равностепенно непрерывным по $s$:
для любых $s<t$ имеем оценку
\[
\bigl|X^{(\alpha)}(t)-X^{(\alpha)}(s)\bigr|
\le
\int_s^t \bigl|F(X^{(\alpha)}(u))\bigr|\,du
+
\bigl|R^{(\alpha)}(s,t)\bigr|,
\]
где первый член пропорционален $|t-s|$,
а второй исчезает при $\alpha\to0$.

Следовательно, по теореме Арцела–Асколи
существует подпоследовательность $\alpha_k\to0$
такая, что $X^{(\alpha_k)}$
сходится равномерно на $[0,T]$
к некоторой непрерывной функции $x(s)$.

Переходя к пределу по этой подпоследовательности
в равенстве для приращений,
и используя непрерывность функции $F$,
получаем, что для любых $0\le s<t\le T$ выполняется
интегральное соотношение
\[
x(t)-x(s)
=
\int_s^t F\bigl(x(u)\bigr)\,du.
\]

Из этого следует, что функция $x$ абсолютно непрерывна
и почти всюду дифференцируема,
причём её производная удовлетворяет
\[
\frac{dx}{ds}=F(x(s))
\quad\text{для почти всех } s\in[0,T].
\]
Непрерывность правой части обеспечивает,
что равенство выполняется для всех $s\in[0,T]$.

Наконец, поскольку обыкновенное дифференциальное уравнение
\[
\frac{dx}{ds}=F(x)
\]
с заданным начальным условием
имеет единственное решение,
предел $x(s)$ не зависит от выбора подпоследовательности.
Следовательно, вся последовательность процессов
$X^{(\alpha)}(s)$
сходится по вероятности,
равномерно по $s\in[0,T]$,
к функции $x(s)$.
\end{proof}


\begin{remark}
Полученная непрерывная динамика описывает усреднённое поведение
разности Q-значений на конечных временных интервалах.
Флуктуации, связанные со случайным членом,
исчезают в этом пределе,
однако они становятся существенными
при анализе поведения системы
на более тонких временных масштабах.
\end{remark}



% -------------------------------------------------
% -------------------------------------------------
% -------------------------------------------------
\subsection{Флуктуации порядка $\sqrt{\alpha}$}

В предыдущем подразделе было доказано,
что при $\alpha \to 0$
интерполированный процесс
\[
X^{(\alpha)}(s)
:=
X_{\lfloor s/\alpha \rfloor}
\]
является равномерно ограниченным и равностепенно непрерывным
на любом конечном интервале $[0,T]$
и, следовательно, по теореме Арцела--Асколи
имеет подпоследовательность,
сходящуюся равномерно к предельной функции $x(s)$.
Было также показано, что любая такая предельная функция
удовлетворяет интегральному уравнению
\[
x(t)-x(s)
=
\int_s^t F\bigl(x(u)\bigr)\,du,
\]
и потому является решением обыкновенного дифференциального уравнения
\[
\frac{dx}{ds}=F(x),
\qquad x(0)=\lim_{\alpha\to0} X^{(\alpha)}(0).
\]

Для изучения случайных отклонений
от детерминированной траектории $x(s)$
введём центрированный и масштабированный процесс
\begin{equation}
\label{eq:def_Y_alpha}
Y^{(\alpha)}(s)
:=
\frac{X^{(\alpha)}(s)-x(s)}{\sqrt{\alpha}}.
\end{equation}

Выбор масштаба $\sqrt{\alpha}$ является естественным,
поскольку за время порядка $1/\alpha$
в динамике накапливается сумма порядка $1/\alpha$
шумовых приращений,
каждое из которых имеет порядок $\alpha$,
а их суммарная дисперсия имеет порядок единицы.

% -------------------------------------------------
\subsection{Диффузионное приближение}

\begin{theorem}[Диффузионный предел для флуктуаций]
\label{thm:sde_limit_correct}
Пусть выполняются следующие условия:
\begin{itemize}
\item функция $F$ непрерывна и локально липшицева;
\item последовательность случайных величин $\eta(m)$ удовлетворяет
\[
\mathbb E[\eta(m)\mid \mathcal F_m]=0;
\]
\item условная дисперсия при фиксированном значении $\Delta Q(m)=x$
имеет вид
\[
\mathbb E[\eta(m)^2\mid \Delta Q(m)=x]
=
\Sigma^2(x),
\]
где функция $\Sigma(x)$ непрерывна;
\item начальные условия удовлетворяют
\[
Y^{(\alpha)}(0)
:=
\frac{X^{(\alpha)}(0)-x(0)}{\sqrt{\alpha}}
\;\xrightarrow{\;\mathbb P\;}\;
Y(0).
\]
\end{itemize}
Тогда при $\alpha\to0$
процесс флуктуаций $Y^{(\alpha)}$
сходится по распределению в пространстве $C([0,T])$
к процессу $Y$,
удовлетворяющему стохастическому дифференциальному уравнению
\[
dY_s
=
F'(x(s))\,Y_s\,ds
+
\Sigma(x(s))\,dW_s,
\qquad
Y(0)\ \text{задано}.
\]
\end{theorem}

\begin{proof}
Дискретная динамика имеет вид
\[
X_{m+1}
=
X_m
+
\alpha F(X_m)
+
\alpha \eta(m).
\]
Рассмотрим кусочно-постоянную интерполяцию
\[
X^{(\alpha)}(s)
=
X_{\lfloor s/\alpha \rfloor}.
\]
Из дискретного уравнения следует представление
\begin{equation}
\label{eq:integral_representation_corrected}
X^{(\alpha)}(s)
=
X^{(\alpha)}(0)
+
\sum_{k=0}^{\lfloor s/\alpha \rfloor-1}
\alpha F(X_k)
+
\sum_{k=0}^{\lfloor s/\alpha \rfloor-1}
\alpha \eta(k).
\end{equation}

Разность между суммой детерминированных членов
и интегралом
\[
\int_0^s F\bigl(X^{(\alpha)}(u)\bigr)\,du
\]
стремится к нулю равномерно по $s\in[0,T]$
при $\alpha\to0$,
поскольку функция $F$ равномерно непрерывна
на компактных множествах,
а разбиение интервала имеет шаг $\alpha$.

Следовательно,
\[
X^{(\alpha)}(s)
=
X^{(\alpha)}(0)
+
\int_0^s F\bigl(X^{(\alpha)}(u)\bigr)\,du
+
M^{(\alpha)}(s)
+
o_{\mathbb P}(1),
\]
где
\[
M^{(\alpha)}(s)
:=
\sum_{k=0}^{\lfloor s/\alpha \rfloor-1}
\alpha \eta(k).
\]

Пусть $x(s)$ — решение детерминированного уравнения
\[
x(s)
=
x(0)
+
\int_0^s F(x(u))\,du.
\]
Определим процесс флуктуаций
\[
Y^{(\alpha)}(s)
:=
\frac{X^{(\alpha)}(s)-x(s)}{\sqrt{\alpha}}.
\]
Тогда
\begin{equation}
\label{eq:Y_decomposition_final_corrected}
Y^{(\alpha)}(s)
=
Y^{(\alpha)}(0)
+
\int_0^s
\frac{F(X^{(\alpha)}(u))-F(x(u))}{\sqrt{\alpha}}\,du
+
\frac{M^{(\alpha)}(s)}{\sqrt{\alpha}}
+
o_{\mathbb P}(1).
\end{equation}

Рассмотрим асимптотику первого интегрального члена
в правой части \eqref{eq:Y_decomposition_final_corrected}.
Напомним, что функция дрейфа $F$
получена как условное математическое ожидание полного приращения
и, вообще говоря, может зависеть от состояния всей системы.
В дальнейшем под $F(x)$ мы понимаем зависимость от первой переменной,
вычисленную вдоль предельной детерминированной траектории $x(\cdot)$.

Предположим, что функция $F$ принадлежит классу $C^1_{\mathrm{loc}}$
по первой переменной.
Так как $X^{(\alpha)} \to x$ по вероятности равномерно на $[0,T]$,
существует компактное множество $K \subset \mathbb R$ такое, что
\[
\mathbb P\!\left(
X^{(\alpha)}(u)\in K,\; x(u)\in K
\ \text{для всех } u\in[0,T]
\right)\xrightarrow{\alpha\to0}1.
\]
На этом компакте производная $F'$ ограничена и равномерно непрерывна.

Для любых $x,y\in K$ по формуле Ньютона--Лейбница имеем тождество
\[
F(y)-F(x)
=
\int_0^1 F'\bigl(x+\theta(y-x)\bigr)\,(y-x)\,d\theta.
\]
Добавляя и вычитая $F'(x)(y-x)$, получаем точное разложение
\[
F(y)-F(x)
=
F'(x)(y-x)
+
(y-x)
\int_0^1
\bigl[
F'\bigl(x+\theta(y-x)\bigr)-F'(x)
\bigr]
\,d\theta.
\]

Применяя это тождество к
$y = X^{(\alpha)}(u)$ и $x = x(u)$,
для каждого $u\in[0,T]$ получаем
\[
F(X^{(\alpha)}(u))-F(x(u))
=
F'(x(u))\,(X^{(\alpha)}(u)-x(u))
+
R^{(\alpha)}(u),
\]
где остаточный член имеет вид
\[
R^{(\alpha)}(u)
=
(X^{(\alpha)}(u)-x(u))
\int_0^1
\bigl[
F'\bigl(x(u)+\theta(X^{(\alpha)}(u)-x(u))\bigr)
-
F'(x(u))
\bigr]
\,d\theta.
\]

Так как
\[
X^{(\alpha)}(u)-x(u)=\sqrt{\alpha}\,Y^{(\alpha)}(u),
\]
то
\[
\frac{R^{(\alpha)}(u)}{\sqrt{\alpha}}
=
Y^{(\alpha)}(u)
\int_0^1
\bigl[
F'\bigl(x(u)+\theta\sqrt{\alpha}Y^{(\alpha)}(u)\bigr)
-
F'(x(u))
\bigr]
\,d\theta.
\]

Из ограниченности по вероятности семейства $Y^{(\alpha)}$
в пространстве $C([0,T])$ следует, что для любого $\varepsilon>0$
существует $M>0$ такое, что
\[
\mathbb P\!\left(
\sup_{u\in[0,T]} |Y^{(\alpha)}(u)| \le M
\right)\ge 1-\varepsilon.
\]
На этом событии аргумент функции $F'$
отличается от $x(u)$ не более чем на $\sqrt{\alpha}M$.
Поскольку $F'$ равномерно непрерывна на компакте $K$,
получаем
\[
\sup_{u\in[0,T]}
\left|
\frac{R^{(\alpha)}(u)}{\sqrt{\alpha}}
\right|
\xrightarrow{\;\mathbb P\;}
0.
\]

Следовательно,
\[
\int_0^s
\frac{F(X^{(\alpha)}(u))-F(x(u))}{\sqrt{\alpha}}
\,du
=
\int_0^s
F'(x(u))\,Y^{(\alpha)}(u)\,du
+
o_{\mathbb P}(1).
\]

Рассмотрим теперь шумовой член
\[
\frac{M^{(\alpha)}(s)}{\sqrt{\alpha}}
=
\sum_{k=0}^{\lfloor s/\alpha \rfloor-1}
\sqrt{\alpha}\,\eta(k).
\]
По предположению о дисперсии при условии
$\Delta Q(k)=x$ имеем
\[
\mathbb E[\eta(k)^2\mid \Delta Q(k)=x]
=
\Sigma^2(x).
\]
Так как $\Delta Q(k)=X^{(\alpha)}(k\alpha)$,
из равномерной сходимости $X^{(\alpha)}\to x$
и непрерывности функции $\Sigma^2$
следует
\[
\mathbb E[\eta(k)^2]
=
\Sigma^2(x(k\alpha))
+
o(1),
\quad
\text{равномерно по } k.
\]

Рассмотрим равномерное разбиение отрезка $[0,s]$
с шагом $\alpha$:
\[
t_k = k\alpha,
\qquad
k=0,1,\dots,\lfloor s/\alpha \rfloor.
\]
Тогда
\[
\sum_{k=0}^{\lfloor s/\alpha \rfloor-1}
\alpha\,\Sigma^2(x(k\alpha))
=
\sum_{k=0}^{\lfloor s/\alpha \rfloor-1}
\Sigma^2(x(t_k))\,(t_{k+1}-t_k),
\]
что является суммой Римана для функции
$u\mapsto \Sigma^2(x(u))$.
Поскольку эта функция непрерывна на $[0,s]$,
получаем
\[
\sum_{k=0}^{\lfloor s/\alpha \rfloor-1}
\alpha\,\Sigma^2(x(k\alpha))
\xrightarrow{\alpha\to0}
\int_0^s \Sigma^2(x(u))\,du.
\]

Рассмотрим фиксированное $s\in[0,T]$.
Случайная величина
\[
\frac{M^{(\alpha)}(s)}{\sqrt{\alpha}}
=
\sum_{k=0}^{\lfloor s/\alpha \rfloor-1}
\sqrt{\alpha}\,\eta(k)
\]
является суммой большого числа центрированных слагаемых
с суммарной дисперсией
\[
\sum_{k=0}^{\lfloor s/\alpha \rfloor-1}
\alpha\,\mathbb E[\eta(k)^2]
\;\xrightarrow{\;}
\int_0^s \Sigma^2(x(u))\,du.
\]
Следовательно, по центральной предельной теореме
\[
\frac{M^{(\alpha)}(s)}{\sqrt{\alpha}}
\;\Longrightarrow\;
\mathcal N\!\left(0,\;\int_0^s \Sigma^2(x(u))\,du\right).
\]

Аналогично, для любых $0\le s_1< s_2\le T$
приращение
\[
\frac{M^{(\alpha)}(s_2)-M^{(\alpha)}(s_1)}{\sqrt{\alpha}}
\]
асимптотически не коррелировано
с $\frac{M^{(\alpha)}(s_1)}{\sqrt{\alpha}}$,
что в пределе приводит к независимости приращений.

Таким образом, конечномерные распределения процесса
$\frac{M^{(\alpha)}}{\sqrt{\alpha}}$
сходятся к гауссовским распределениям
с ковариационной функцией
\[
\mathbb E[Z(s)Z(t)]
=
\int_0^{\min(s,t)} \Sigma^2(x(u))\,du.
\]
Поэтому предельный процесс может быть представлен
в виде стохастического интеграла
\[
Z(s)=\int_0^s \Sigma(x(u))\,dW_u,
\]
где $W$ — стандартный винеровский процесс.


Переходя к пределу в \eqref{eq:Y_decomposition_final_corrected},
получаем, что любой предельный процесс $Y$
удовлетворяет интегральному уравнению
\[
Y(s)
=
Y(0)
+
\int_0^s F'(x(u))\,Y(u)\,du
+
\int_0^s \Sigma(x(u))\,dW_u.
\]
Линейность этого уравнения обеспечивает
единственность решения,
что завершает доказательство.

\end{proof}


% -------------------------------------------------
\subsection{Локальная линейная аппроксимация вблизи дефектного равновесия}

Дальнейший анализ посвящён поведению процесса флуктуаций
в окрестности дефектного равновесия $x^\star$,
определённого как устойчивое равновесие
детерминированной предельной динамики,
то есть
\[
F(x^\star)=0,
\qquad
F'(x^\star)<0.
\]

Напомним, что процесс флуктуаций $Y$
удовлетворяет стохастическому дифференциальному уравнению
\[
dY_s
=
F'(x(s))\,Y_s\,ds
+
\Sigma(x(s))\,dW_s,
\]
где $x(s)$ — решение детерминированного уравнения
\[
\dot x = F(x),
\qquad
x(0)\ \text{задано}.
\]

\begin{theorem}[Локальная линейная аппроксимация]
\label{thm:ou_correct}
Существует $\delta>0$ такое, что
если $|x(0)-x^\star|<\delta$,
то на временных интервалах,
пока $x(s)$ остаётся в $\delta$-окрестности точки $x^\star$,
процесс флуктуаций $Y$
аппроксимируется решением линейного стохастического уравнения
\[
dZ_s
=
-\kappa\,Z_s\,ds
+
\Sigma(x^\star)\,dW_s,
\qquad
\kappa := -F'(x^\star)>0,
\]
то есть процессом Орнштейна--Уленбека.
\end{theorem}

\begin{proof}
Поскольку $F\in C^1_{\mathrm{loc}}$ по первой переменной,
из условий $F(x^\star)=0$ и $F'(x^\star)<0$ следует,
что в окрестности $x^\star$ справедливо разложение
\[
F'(x)
=
F'(x^\star)
+
o(1)
\quad
\text{при } x\to x^\star.
\]

Если начальное условие $x(0)$ достаточно близко к $x^\star$,
то устойчивость равновесия гарантирует,
что $x(s)$ остаётся в малой окрестности $x^\star$
на рассматриваемом временном интервале.
Следовательно,
\[
F'(x(s)) = -\kappa + o(1),
\qquad
\Sigma(x(s)) = \Sigma(x^\star) + o(1),
\]
где остаточные члены стремятся к нулю равномерно по $s$
на каждом фиксированном конечном интервале времени.

Отбрасывая высшие порядки малости,
получаем линейное стохастическое уравнение
с постоянными коэффициентами,
которое и совпадает с уравнением Орнштейна--Уленбека.
\end{proof}

\begin{remark}
Полученная линейная аппроксимация описывает
типичное поведение флуктуаций
вблизи устойчивого дефектного равновесия.
Именно эта аппроксимация будет использована
в следующем разделе
для анализа времени выхода процесса
из малой окрестности точки $x^\star$.
\end{remark}


% =================================================
% =================================================
\section{Ловушка дефекции и время выхода}

В данном разделе мы формализуем понятие ловушки дефекции
и выводим асимптотическую формулу
для среднего времени выхода из неё
в пределе малого шага обучения $\alpha \to 0$.

% -------------------------------------------------
\subsection{Определение ловушки дефекции}

\begin{definition}[Дефектная ловушка]
\label{def:trap}
Для фиксированного $\varepsilon \in (0,1)$
дефектная ловушка определяется как множество
\[
\mathcal{T}_\varepsilon
=
\{ x \in \mathbb{R} : \sigma(\beta x) \le \varepsilon \}.
\]
Граница ловушки имеет координату
\[
x_{\mathrm{bd}}
=
\frac{1}{\beta}
\log\frac{\varepsilon}{1-\varepsilon}.
\]
Расстояние от дефектного равновесия $x^\star$
до границы ловушки равно
\[
\Delta := x_{\mathrm{bd}} - x^\star > 0.
\]
\end{definition}

\begin{remark}
Множество $\mathcal{T}_\varepsilon$
соответствует области состояний,
в которых вероятность кооперации
экспоненциально мала.
\end{remark}

% -------------------------------------------------
\subsection{Формулировка задачи о времени выхода}

Вблизи дефектного равновесия $x^\star$
динамика исходного процесса обучения
может быть аппроксимирована
одномерным стохастическим процессом
с малым шумом.
В рамках локальной линейной аппроксимации
(см. Теорему~\ref{thm:ou_correct})
мы рассматриваем процесс $X_s$,
удовлетворяющий стохастическому уравнению
\[
dX_s
=
-\kappa (X_s - x^\star)\,ds
+
\sqrt{\alpha}\,\Sigma(x^\star)\,dW_s,
\]
где $\kappa>0$,
а параметр $\alpha$ характеризует величину шума.

\begin{definition}[Время выхода]
\label{def:exit_time}
Временем выхода из дефектной ловушки
называется случайная величина
\[
\tau_{\mathrm{exit}}
=
\inf\{ s \ge 0 : X_s \notin \mathcal{T}_\varepsilon \}.
\]
\end{definition}

Наша цель состоит в исследовании асимптотики
математического ожидания
$\mathbb{E}[\tau_{\mathrm{exit}}]$
при $\alpha \to 0$.


% -------------------------------------------------
\subsection{Квазипотенциал и энергетический барьер}

Рассмотрим стохастический процесс
\[
dX_s
=
-\kappa (X_s - x^\star)\,ds
+
\sqrt{\alpha}\,\Sigma(x^\star)\,dW_s,
\]
который используется в качестве локальной модели
динамики исходного процесса обучения
вблизи дефектного равновесия $x^\star$.

Детерминированная часть этого уравнения
имеет вид
\[
-\kappa (x - x^\star)
=
- \frac{d}{dx} V(x),
\]
где функция
\[
V(x)
=
\frac{\kappa}{2}(x - x^\star)^2
\]
играет роль потенциальной энергии.
В этом смысле процесс $X_s$
является одномерной градиентной диффузией
с малым шумом.

В теории больших уклонений
для стохастических дифференциальных уравнений
с малым шумом
ключевую роль играет \emph{квазипотенциал},
который описывает асимптотику вероятностей
редких переходов между состояниями.
В случае одномерной градиентной диффузии
квазипотенциал между точками
совпадает с разностью значений функции $V$.

\begin{lemma}[Высота потенциального барьера]
\label{lem:barrier}
Разность квазипотенциала
между устойчивым равновесием $x^\star$
и границей дефектной ловушки $x_{\mathrm{bd}}$
равна
\[
\Delta V
=
V(x_{\mathrm{bd}}) - V(x^\star)
=
\frac{\kappa}{2}\,(x_{\mathrm{bd}} - x^\star)^2
=
\frac{\kappa \Delta^2}{2}.
\]
\end{lemma}

\begin{proof}
По определению функции $V$
имеем $V(x^\star)=0$.
Подставляя $x=x_{\mathrm{bd}}$
и используя обозначение
$\Delta = x_{\mathrm{bd}} - x^\star$,
получаем
\[
\Delta V
=
\frac{\kappa}{2}\Delta^2,
\]
что и требовалось доказать.
\end{proof}

\begin{remark}
Величина $\Delta V$
характеризует энергетический барьер,
который процесс должен преодолеть
для выхода из дефектной ловушки.
В следующем разделе будет показано,
что именно эта величина
определяет экспоненциальный масштаб
среднего времени выхода
при $\alpha \to 0$.
\end{remark}


% -------------------------------------------------
\subsection{Теорема о времени выхода}

\begin{theorem}[Экспоненциальная асимптотика времени выхода]
\label{thm:exit_time}
Рассмотрим процесс $X_s$,
удовлетворяющий стохастическому уравнению
\[
dX_s
=
-\kappa (X_s - x^\star)\,ds
+
\sqrt{\alpha}\,\Sigma(x^\star)\,dW_s,
\]
и пусть начальное условие $X_0$
принадлежит фиксированной окрестности точки $x^\star$,
не зависящей от $\alpha$.
Тогда при $\alpha \to 0$
математическое ожидание времени выхода
$\tau_{\mathrm{exit}}$
из дефектной ловушки $\mathcal T_\varepsilon$
обладает логарифмической асимптотикой
\[
\lim_{\alpha\to0}
\alpha \log \mathbb E[\tau_{\mathrm{exit}}]
=
\frac{\kappa \Delta^2}{\Sigma^2(x^\star)}.
\]
В частности,
\[
\mathbb E[\tau_{\mathrm{exit}}]
=
\exp\!\left(
\frac{\kappa \Delta^2}
{\alpha \Sigma^2(x^\star)}
\right)
\cdot
(1+o(1)).
\]
\end{theorem}

\begin{proof}
Рассматриваемый процесс является одномерной диффузией
с малым шумом порядка $\sqrt{\alpha}$
и единственным устойчивым равновесием $x^\star$
внутри области $\mathcal T_\varepsilon$.
Для таких процессов теория больших уклонений
для стохастических дифференциальных уравнений
с малым шумом
(см.\ теорию Фрейдлина--Вентцеля)
утверждает,
что логарифмическая асимптотика
среднего времени выхода
из области притяжения устойчивого равновесия
определяется высотой квазипотенциального барьера:
\[
\lim_{\alpha\to0}
\alpha \log \mathbb E[\tau_{\mathrm{exit}}]
=
\frac{2\Delta V}{\Sigma^2(x^\star)}.
\]

В рассматриваемом случае,
согласно Лемме~\ref{lem:barrier},
разность квазипотенциала равна
\[
\Delta V
=
\frac{\kappa \Delta^2}{2}.
\]
Подстановка этого выражения
даёт
\[
\lim_{\alpha\to0}
\alpha \log \mathbb E[\tau_{\mathrm{exit}}]
=
\frac{\kappa \Delta^2}{\Sigma^2(x^\star)}.
\]

Утверждение теоремы
о форме экспоненциальной асимптотики
следует из данного равенства.
Предэкспоненциальный множитель
зависит от локальной геометрии коэффициентов
и в данной работе не анализируется,
поскольку он не влияет
на главный экспоненциальный порядок.
\end{proof}


% -------------------------------------------------
\subsection{Интерпретация результата}

\begin{remark}
Экспоненциальный характер времени выхода
обусловлен совместным действием следующих факторов:
\begin{itemize}
\item малой интенсивности стохастических флуктуаций,
пропорциональной $\sqrt{\alpha}$;
\item линейной силы возврата к дефектному равновесию,
характеризуемой коэффициентом $\kappa$;
\item конечного расстояния $\Delta$
между дефектным равновесием и границей ловушки.
\end{itemize}
В результате даже умеренное уменьшение шага обучения $\alpha$
приводит к резкому увеличению
типичного времени пребывания системы
в дефектной области.
\end{remark}

\begin{remark}
Полученная асимптотика
даёт теоретическое объяснение
наблюдаемому в численных экспериментах эффекту
длительного удержания системы
в режиме дефекции
при малых шагах обучения.
В рамках построенной модели
такое поведение является следствием
редкости флуктуаций,
достаточных для выхода из дефектной ловушки,
а не наличия дополнительных поглощающих состояний.
\end{remark}



% =================================================
\section{Заключение}

В работе построена строгая асимптотическая теория
дефектных ловушек
в многоагентном Q-обучении
в пределе малого шага обучения.
Исходная дискретная динамика обучения
была сведена к одномерному стохастическому процессу,
для которого последовательно получены
детерминированный предел,
диффузионное приближение
и локальная линейная аппроксимация
вблизи дефектного равновесия.

Показано, что в режиме малого шума
выход из дефектной области
является редким событием,
а среднее время выхода
обладает экспоненциальной асимптотикой
по параметру $\alpha^{-1}$.
Главный экспоненциальный порядок
определяется высотой квазипотенциального барьера
и не зависит от деталей предэкспоненциального множителя.

Полученные результаты
объясняют устойчивость дефектного поведения,
наблюдаемую в численных экспериментах
по многоагентному обучению с подкреплением,
и демонстрируют,
что она может возникать
как следствие динамики обучения,
даже при отсутствии явно поглощающих состояний.

\end{document}
